% huerden-baum.tex — semio TikZ example: a branch of the Hürden taxonomy.
% Compile:  python _scripts/render_tikz.py huerden-baum
% Embed:    inside a \begin{Figure}[title=...]\SemioImage[fit=none]{...huerden-baum.pdf}\end{Figure}
\documentclass[border=6pt]{standalone}
\usepackage{semio-tikz}
\begin{document}
\begin{tikzpicture}[node distance=9mm]

  % root + one domain branch (H) expanded to a family and two leaves
  \node[semio hub, text width=52mm] (root)
    {\SemioNode{Hürdensystem}{Digitale Wiederverwendungs-Hürden}};
  \node[semio hub, text width=52mm, below=of root] (h)
    {\SemioNode{H}{Physische Rückgewinnung \& Ausführung}};
  \node[semio hub, text width=40mm, below=of h] (h5)
    {\SemioNode{H5}{Lagerfähigkeit}};
  \node[semio node, text width=34mm, below left=14mm and -6mm of h5] (h51)
    {\SemioNode{H5.1}{Keine Lagerkapazität}};
  \node[semio node, text width=34mm, below right=14mm and -6mm of h5] (h53)
    {\SemioNode{H5.3}{Ungeeignete Lagerbedingungen}};

  \draw[semio edge] (root) -- (h);
  \draw[semio edge] (h) -- (h5);
  \draw[semio edge] (h5) -- (h51);
  \draw[semio edge] (h5) -- (h53);

  % register / legend
  \node[semio register, text width=44mm, right=16mm of h, anchor=west] (leg)
    {HAT\_UNTERKATEGORIE ordnet 409 Knoten\\
     in 8 Bereiche A–H.\\[2pt]
     {\semiomonofont barriere\_code} als Chip,\\
     Kategorien heller, Blätter neutral.};

\end{tikzpicture}
\end{document}
